% the abstract
\begin{abstractzh}
在大模型自进化背景下，训练任务的分布可能不均衡，新数据难以保证分布均匀，直接参与强化学习或微调会导致模型此前掌握的能力有一定程度退化。本文面向该问题，提出一套面向智能体大模型的持续强化学习方法，核心包括：（1）复合持续学习目标，将 GRPO 策略梯度与逆向 KL 锚定、经验回放、熵正则相融合；（2）由外部冻结大模型裁判对轨迹的完成度、正确性、轨迹质量、安全四个维度打分，并以二进制安全门控乘性聚合；（3）按能力域划分的九桶经验回放缓冲区，以平方根加权分配容量、桶内独立淘汰、设计能力域的多维空间，按空间距离选桶回放，防止不同能力互相侵占；（4）双模型任务难度标注与按桶顺序训练。方法以零框架改动的方式注入 verl 强化学习框架。实验在 Qwen3.5-9B 基座与 ClawEval 纯文本子集（195 个任务、9 个能力域）上开展，采用训完统一评测的抗遗忘评测协议，与无持续学习机制的同规模基线对照。实验结果表明，本方法在按桶顺序训练下能有效缓解早期训练能力的遗忘。
\end{abstractzh}